% Section 6.3, from lectures/L06/appendix/b_exercises.md.

\section{Exercises}
\label{c6:sec:exercises}\label{c6:app:b}

\begin{aside}
\textbf{How to work these.} These exercises reinforce the concepts from \secref{c6:sec:threads}.
Each set names the file or directory to work in: write it in its own directory under
\file{lectures/L06/exercises} in the companion repository, with the file names its README gives.
Build it with the Makefile from \secref{c1:sec:makefile}, and add \code{-pthread} to the compiler
flags (\code{CXX_FLAGS := -Wall -Werror -std=c++17 -pthread}). GCC requires it for programs that
use \code{std::thread}; without it, the program may fail to link, or fail at runtime, on systems
with a C library older than glibc 2.34.

\textbf{How to check your work.} Run \code{make test} from the repository root. A set's tests
switch on as soon as its files exist; a set you have not started is reported as \code{SKIP}. Each
file grows through its set, and the tests check where it ends up. Threads run in whatever order the
scheduler picks, so the tests check what a correct program guarantees, such as exact counts and
threads that stop when asked, rather than the order of its output.

The solutions are published in the companion repository, one program per set under
\file{lectures/L06/appendix/solutions} (Sets~4 and~5 share one), with the answers to the reflection
questions in the README beside them. Try each exercise yourself before looking at them.
\end{aside}

% ------------------------------------------------------------------------------------------
\exerciseset{Threads and Atomic Stop Flag}
\label{c6:set:1}
Create a file named \file{exercise1.cpp} to work in for this exercise set.

\exercise{Creating a worker thread}{Code}
\label{c6:ex:1.1}
In this exercise you will create a simple worker function that runs in a separate thread.

\task{Tasks}
In \file{exercise1.cpp}, create a function named \code{workerThread()}.

The function shall:
\begin{itemize}
  \item Take a \code{std::uint8_t} named \code{printCount}: The number of times to print.
  \item Take a \code{std::uint16_t} named \code{printSpeed_ms}: The delay in milliseconds between
    each print.
  \item Return \code{void}.
  \item Print \code{"Worker thread running!"} \code{printCount} times.
  \item Sleep for \code{printSpeed_ms} milliseconds between each print using
    \code{std::this_thread::sleep_for()} with \code{std::chrono::milliseconds}.
\end{itemize}

\exercise{Starting and joining the thread}{Code}
\label{c6:ex:1.2}
In \file{exercise1.cpp}:
\begin{itemize}
  \item Create a \code{std::thread} that runs \code{workerThread()}, passing \code{5U} as
    \code{printCount} and \code{500U} as \code{printSpeed_ms}.
  \item Print \code{"Main thread running!"} from \code{main()}.
  \item Call \code{join()} on the worker thread before returning from \code{main()}.
\end{itemize}

\task{Reflection}
\begin{itemize}
  \item What happens if \code{join()} is not called?
  \item Which thread prints first, the main thread or the worker thread?
  \item Is the print order guaranteed?
\end{itemize}

\exercise{Atomic stop flag}{Code}
\label{c6:ex:1.3}
In this exercise you will modify the worker thread so that it runs until a stop flag is set.

\task{Tasks}
Create an atomic boolean variable in \code{main()}:

\begin{cppcode}
std::atomic<bool> stop{false};
\end{cppcode}

\noindent Modify \code{workerThread()} so that it:
\begin{itemize}
  \item Replaces the \code{printCount} parameter with a \code{const std::atomic<bool>&} named
    \code{stop}, placed after \code{printSpeed_ms}.
  \item Runs as long as \code{stop.load()} is \code{false}.
  \item Prints \code{"Worker thread running!"} once every \code{printSpeed_ms} milliseconds.
\end{itemize}
In \code{main()}:
\begin{itemize}
  \item Start the worker thread.
  \item Let the main thread sleep for \code{3000 ms}.
  \item Set the stop flag using \code{stop.store(true)}.
  \item Join the worker thread.
\end{itemize}

\task{Reflection}
\begin{itemize}
  \item Why is \code{std::atomic<bool>} suitable for this stop flag?
  \item Would a normal \code{bool} be safe here?
  \item Why is \code{load()} used when reading the flag?
  \item Why is \code{store()} used when writing the flag?
\end{itemize}

% ------------------------------------------------------------------------------------------
\exerciseset{Shared Counter with Data Race}
\label{c6:set:2}
Create a file named \file{exercise2.cpp} to work in for this exercise set.

\exercise{Shared counter without synchronization}{Code}
\label{c6:ex:2.1}
In this exercise you will intentionally create a data race.

\task{Tasks}
Create a shared counter variable in \code{main()}:

\begin{cppcode}
std::uint32_t counter{};
\end{cppcode}

\noindent Create a function named \code{incrementCounter()}. The function shall:
\begin{itemize}
  \item Take a reference to the counter.
  \item Take a number of iterations (\code{std::uint32_t}).
  \item Increment the counter once per iteration.
  \item Be marked \code{noexcept}.
\end{itemize}
In \code{main()}:
\begin{itemize}
  \item Create two threads that both call \code{incrementCounter()}.
  \item Let each thread increment the counter \code{100000U} times.
  \item Join both threads.
  \item Print the final counter value.
\end{itemize}

\task{Reflection}
\begin{itemize}
  \item What final value do you expect?
  \item Do you always get the expected value?
  \item Why is this program unsafe?
  \item Which condition makes this a data race?
\end{itemize}

\exercise{Protecting the counter with a mutex}{Code}
\label{c6:ex:2.2}
In this exercise you will fix the data race using \code{std::mutex}.

\task{Tasks}
Create a mutex in \code{main()}:

\begin{cppcode}
std::mutex mutex{};
\end{cppcode}

\noindent Modify \code{incrementCounter()} so that it:
\begin{itemize}
  \item Takes a reference to the mutex.
  \item Locks the mutex before incrementing the counter.
  \item Uses \code{std::lock_guard<std::mutex>}.
  \item Unlocks automatically when the lock guard goes out of scope.
\end{itemize}

\task{Reflection}
\begin{itemize}
  \item What final counter value do you get now?
  \item Why does the mutex solve the problem?
  \item Why is \code{std::lock_guard} safer than manually calling \code{lock()} and
    \code{unlock()}?
\end{itemize}

\exercise{Atomic counter}{Code}
\label{c6:ex:2.3}
In this exercise you will replace the mutex-protected counter with an atomic counter.

\task{Tasks}
Replace the counter with:

\begin{cppcode}
std::atomic<std::uint32_t> counter{};
\end{cppcode}

\noindent Modify \code{incrementCounter()} so that it:
\begin{itemize}
  \item Takes a reference to the atomic counter.
  \item Increments it once per iteration.
\end{itemize}

\task{Reflection}
\begin{itemize}
  \item What final counter value do you get?
  \item Why is \code{std::atomic<std::uint32_t>} sufficient in this exercise?
  \item Would an atomic counter still be sufficient if several related variables had to be updated
    together?
\end{itemize}

% ------------------------------------------------------------------------------------------
\exerciseset{Shared Data Protected by Mutex}
\label{c6:set:3}
Create a file named \file{exercise3.cpp} to work in for this exercise set.

\exercise{Shared memory structure}{Code}
\label{c6:ex:3.1}
In this exercise you will create shared data that must be protected by a mutex.

\task{Tasks}
In \file{exercise3.cpp}, create a structure named \code{SharedMem} with the following members:
\begin{itemize}
  \item \code{data}: A 16-bit unsigned value.
  \item \code{newData}: True if new data is available, false otherwise.
\end{itemize}
The structure represents data shared between a transmitter thread and a receiver thread.

\exercise{Transmitter thread}{Code}
\label{c6:ex:3.2}
In \file{exercise3.cpp}, create a function named \code{txThread()}.

The function shall:
\begin{itemize}
  \item Take a reference to \code{SharedMem}.
  \item Take a reference to \code{std::mutex}.
  \item Take a reference to \code{const std::atomic<bool>} named \code{stop}.
  \item Run until the stop flag is set.
  \item Produce an incrementing \code{std::uint16_t} value.
  \item Store the value in \code{shared.data}.
  \item Set \code{shared.newData} to \code{true}.
  \item Sleep for \code{1000 ms} between transmissions.
\end{itemize}
The update of \code{shared.data} and \code{shared.newData} shall be protected using a lock guard.

\exercise{Receiver thread}{Code}
\label{c6:ex:3.3}
In \file{exercise3.cpp}, create a function named \code{rxThread()}.

The function shall:
\begin{itemize}
  \item Take a reference to \code{SharedMem}.
  \item Take a reference to \code{std::mutex}.
  \item Take a reference to \code{const std::atomic<bool>} named \code{stop}.
  \item Run until the stop flag is set.
  \item Check whether new data is available.
  \item Print the received data if \code{shared.newData} is \code{true}.
  \item Set \code{shared.newData} to \code{false} after consuming the data.
  \item Sleep for \code{100 ms} between receive attempts.
\end{itemize}
The check and update of \code{shared.data} and \code{shared.newData} shall be protected by the same
mutex.

\exercise{Running TX and RX threads}{Code}
\label{c6:ex:3.4}
In \file{exercise3.cpp}:
\begin{itemize}
  \item Create one \code{SharedMem} object.
  \item Create one \code{std::mutex}.
  \item Create one atomic stop flag.
  \item Start one TX thread.
  \item Start one RX thread.
  \item Let the main thread sleep for \code{10000 ms}.
  \item Set the stop flag.
  \item Join both threads.
\end{itemize}

\task{Reflection}
\begin{itemize}
  \item Why must \code{data} and \code{newData} be protected by the same mutex?
  \item Why is \code{newData} not atomic in this design?
  \item What could happen if TX updated \code{data} and \code{newData} without locking the mutex?
  \item Why is the stop flag atomic instead of mutex-protected?
  \item Could \code{std::atomic<SharedMem>} be used instead of a mutex? Why or why not?
\end{itemize}

% ------------------------------------------------------------------------------------------
\exerciseset{Thread-Safe Stub Driver}
\label{c6:set:4}
Create the following directory structure to work in for this exercise set:

\begin{console}
Makefile
include/
    driver/
        counter/
            interface.hpp
            stub.hpp
source/
    main.cpp
\end{console}

\note When a class overrides a method that is marked \code{[[nodiscard]]} in the interface, repeat
the attribute on the overriding method as well. The attribute is not inherited; see
\secref{c3:app:b}.

\exercise{Counter driver interface}{Code}
\label{c6:ex:4.1}
In this exercise you will design an interface named \code{driver::counter::Interface}.

The interface represents a generic counter driver.

\task{Tasks}
Design a class named \code{driver::counter::Interface}.

All methods except the destructor shall be pure virtual (\code{= 0}).

\task{a) Destructor}
Add a destructor that is:
\begin{itemize}
  \item Virtual.
  \item Marked \code{noexcept}.
  \item Implemented using \code{= default}.
\end{itemize}

\task{b) Checking initialization state}
Add a pure virtual method \code{isInitialized()} that:
\begin{itemize}
  \item Returns \code{true} if the driver is initialized, \code{false} otherwise.
  \item Does not modify the object.
  \item Is marked \code{noexcept} and \code{[[nodiscard]]}.
\end{itemize}

\task{c) Reading the counter}
Add a pure virtual method \code{value()} that:
\begin{itemize}
  \item Returns the current counter value as \code{std::uint32_t}.
  \item Does not modify the object.
  \item Is marked \code{noexcept} and \code{[[nodiscard]]}.
\end{itemize}

\task{d) Incrementing the counter}
Add a pure virtual method \code{increment()} that:
\begin{itemize}
  \item Increments the counter by one.
  \item Does not return a value.
  \item Is marked \code{noexcept}.
\end{itemize}

\task{e) Resetting the counter}
Add a pure virtual method \code{reset()} that:
\begin{itemize}
  \item Resets the counter to zero.
  \item Does not return a value.
  \item Is marked \code{noexcept}.
\end{itemize}

\exercise{Thread-safe stub counter}{Code}
\label{c6:ex:4.2}
In this exercise you will implement \code{driver::counter::Stub}.

The class shall:
\begin{itemize}
  \item Inherit from \code{driver::counter::Interface}.
  \item Be marked \code{final}.
  \item Be safe to use from multiple threads (via a mutex).
\end{itemize}

\task{a) Member variables}
Add two private member variables:
\begin{itemize}
  \item \code{myMutex}:
  \begin{itemize}
    \item Type: \code{std::mutex}.
    \item Protects access to \code{myValue} below.
    \item Must be marked \code{mutable} because \code{value()} is a \code{const} method but still
      needs to lock the mutex.
  \end{itemize}
  \item \code{myValue}:
  \begin{itemize}
    \item Type: \code{std::uint32_t}.
    \item Stores the counter value.
  \end{itemize}
\end{itemize}

\task{b) Constructor}
Add a default constructor that:
\begin{itemize}
  \item Initializes the counter to zero.
  \item Is marked \code{noexcept}.
\end{itemize}

\task{c) Destructor}
Add a destructor that is:
\begin{itemize}
  \item Public.
  \item Marked \code{noexcept}.
  \item Declared with \code{override}.
  \item Implemented using \code{= default}.
\end{itemize}

\task{d) Implement interface methods}
Implement all methods required by the interface:
\begin{itemize}
  \item \code{isInitialized()} shall return \code{true}.
  \item \code{value()} shall:
  \begin{itemize}
    \item Lock the mutex using \code{std::lock_guard<std::mutex>}.
    \item Return \code{myValue}.
  \end{itemize}
  \item \code{increment()} shall:
  \begin{itemize}
    \item Lock the mutex using \code{std::lock_guard<std::mutex>}.
    \item Increment \code{myValue}.
  \end{itemize}
  \item \code{reset()} shall:
  \begin{itemize}
    \item Lock the mutex using \code{std::lock_guard<std::mutex>}.
    \item Set \code{myValue} to zero.
  \end{itemize}
\end{itemize}

\task{e) Disable copy and move semantics}
Delete the following functions in the public section of the class:
\begin{itemize}
  \item Copy constructor.
  \item Move constructor.
  \item Copy assignment operator.
  \item Move assignment operator.
\end{itemize}

\exercise{Testing the thread-safe stub}{Code}
\label{c6:ex:4.3}
Create a function named \code{counterThread()}.

The function shall:
\begin{itemize}
  \item Take a reference to \code{driver::counter::Interface}.
  \item Take a \code{std::uint32_t} named \code{iterations}, which is the number of times to call
    \code{increment()}.
  \item Be marked \code{noexcept}.
\end{itemize}
In \file{main.cpp}:
\begin{itemize}
  \item Create one \code{driver::counter::Stub}.
  \item Start two threads, both running \code{counterThread()} --- the first with \code{100U}
    iterations, the second with \code{200U}.
  \item Join both threads.
  \item Print the final counter value using \code{value()}.
\end{itemize}

\task{Reflection}
\begin{itemize}
  \item Why does \code{value()} need to lock the mutex even though it only reads the value?
  \item Why is \code{myMutex} marked \code{mutable}?
  \item What would happen if \code{increment()} did not lock the mutex?
  \item Could \code{std::atomic<std::uint32_t>} be used instead in this specific class?
  \item Why might a mutex still be preferable in a more complex driver?
\end{itemize}

% ------------------------------------------------------------------------------------------
\exerciseset{Atomic State and Mutex-Protected Data}
\label{c6:set:5}

\exercise{Driver state}{Code}
\label{c6:ex:5.1}
In this exercise you will extend the counter stub with an initialization state.

\task{Tasks}
Add a private member variable to \code{driver::counter::Stub}:

\begin{cppcode}
std::atomic<bool> myInitialized;
\end{cppcode}

\noindent Initialize it to \code{true} in the constructor.

\exercise{Initialization methods}{Code}
\label{c6:ex:5.2}
Add a public method \code{setInitialized()} specific to the stub. The method shall:
\begin{itemize}
  \item Take a \code{bool} named \code{initialized}.
  \item Write to \code{myInitialized} using \code{store()}.
  \item Be marked \code{noexcept}.
\end{itemize}
Override \code{isInitialized()} from the interface so that it:
\begin{itemize}
  \item Reads \code{myInitialized} using \code{load()}.
  \item Returns the current initialization state instead of always returning \code{true}.
\end{itemize}

\exercise{Guarding counter operations}{Code}
\label{c6:ex:5.3}
Modify \code{increment()}, \code{value()}, and \code{reset()}:
\begin{itemize}
  \item \code{increment()} shall do nothing if the driver is not initialized.
  \item \code{value()} shall return \code{0U} if the driver is not initialized.
  \item \code{reset()} shall do nothing if the driver is not initialized.
\end{itemize}
The counter value itself shall still be protected by the mutex.

\task{Reflection}
\begin{itemize}
  \item Why is \code{myInitialized} suitable as an atomic variable?
  \item Why is \code{myValue} still protected by a mutex?
  \item Would it be safe to make both variables atomic if future requirements added more related
    state?
  \item What is the main design difference between a thread-safe flag and thread-safe shared data?
  \item What could happen if \code{setInitialized(false)} is called between the initialization
    check and the mutex lock in \code{increment()}? How would you fix this?
\end{itemize}

% ------------------------------------------------------------------------------------------
\exerciseset{Condition Variable}
\label{c6:set:6}
Create a file named \file{exercise6.cpp} by copying \file{exercise3.cpp}. All modifications in
this exercise set are made to that file.

\exercise{Replacing polling with a condition variable}{Code}
\label{c6:ex:6.1}
In Exercise~\ref{c6:ex:3.3} the receiver thread checks for new data by polling in a loop with a
\code{100 ms} sleep. In this exercise you will replace that pattern with a
\code{std::condition_variable} so the receiver wakes up immediately when data is available.

\task{Tasks}
Add a \code{std::condition_variable} named \code{cv} alongside the existing mutex.

Create a helper function named \code{hasNewData()} to use as the wake-up predicate. The function
shall:
\begin{itemize}
  \item Take a \code{const SharedMem&}.
  \item Take a \code{const std::atomic<bool>&} named \code{stop}.
  \item Return \code{true} if \code{shared.newData} is \code{true} or the stop flag is set,
    \code{false} otherwise.
  \item Be marked \code{noexcept} and \code{[[nodiscard]]}.
\end{itemize}
Modify \code{txThread()} so that it calls \code{cv.notify_one()} after releasing the lock guard that
protects the data update.

Modify \code{rxThread()} so that instead of sleeping it:
\begin{itemize}
  \item Acquires a \code{std::unique_lock<std::mutex>}.
  \item Calls \code{cv.wait()} with the lock and a predicate that calls \code{hasNewData()}.
\end{itemize}
\note The predicate is passed as a lambda; the \code{[&]} syntax captures \code{shared} and
\code{stop} by reference so the predicate can access both. Lambdas are covered in
\secref{c6:sec:condvar}.

In \code{main()}, set the stop flag while holding the mutex (in a scope of its own with a
\code{std::lock_guard<std::mutex>}), then call \code{cv.notify_all()} to wake the receiver thread
so it can exit cleanly.

\note The stop flag is atomic, but it is also part of the predicate, and a variable the predicate
depends on must be modified while holding the mutex. Otherwise the flag could be set, and
\code{notify_all()} called, just after the receiver found the predicate \code{false} but before it
went to sleep. That notification would be lost, and the receiver would never wake up.

\task{Reflection}
\begin{itemize}
  \item Why is \code{std::unique_lock} required here instead of \code{std::lock_guard}?
  \item What is a spurious wake-up, and how does the predicate protect against it?
  \item Why must the stop flag also be checked inside the predicate?
  \item How does this design differ from the polling approach in Exercise~\ref{c6:ex:3.3}?
  \item What would happen if \code{notify_all()} was not called after setting the stop flag?
  \item Why is \code{notify_all()} used here instead of \code{notify_one()}?
\end{itemize}

% ------------------------------------------------------------------------------------------
\exerciseset{\code{std::async} and \code{std::future}}
\label{c6:set:7}
Create a file named \file{exercise7.cpp} to work in for this exercise set.

\exercise{Offloading a computation}{Code}
\label{c6:ex:7.1}
In this exercise you will use \code{std::async} to run a validation in the background while the
main thread continues working.

\task{Tasks}
Create a function named \code{validateFirmware()}. The function shall:
\begin{itemize}
  \item Take a \code{const std::uint8_t*} named \code{buf} and a \code{std::uint16_t} named
    \code{bufLen}.
  \item Return \code{false} if \code{buf} is \code{nullptr} or \code{bufLen} is \code{0}.
  \item Sleep for \code{2000 ms} to simulate a slow flash read before validating.
  \item Iterate over all bytes in \code{buf}.
  \item Return \code{false} if any byte has the value \code{0xFFU} (unprogrammed flash).
  \item Return \code{true} if the image is valid.
  \item Be marked \code{noexcept} and \code{[[nodiscard]]}.
\end{itemize}
In \code{main()}:
\begin{itemize}
  \item Create a buffer with some test data.
  \item Launch \code{validateFirmware()} asynchronously using \code{std::async} with
    \code{std::launch::async}, storing the result in a \code{std::future<bool>}.
  \item Print \code{"Booting system while validating firmware..."} while validation is running.
  \item Retrieve the result using \code{future.get()} and print whether the firmware is valid or
    not.
\end{itemize}

\task{Reflection}
\begin{itemize}
  \item What does \code{std::launch::async} guarantee compared to omitting the launch policy?
  \item At what point does the main thread block?
  \item What would happen if \code{future.get()} was called immediately after \code{std::async}?
\end{itemize}

\exercise{Checking whether the result is ready}{Code}
\label{c6:ex:7.2}
Modify \code{main()} so that instead of calling \code{future.get()} directly, it polls in a loop:
\begin{itemize}
  \item Each iteration shall call \code{future.wait_for()} with
    \code{std::chrono::milliseconds(200U)} and compare the result against
    \code{std::future_status::ready}.
  \item As long as the result is not ready, print \code{"Still waiting for firmware validation..."}.
  \item Once the loop exits, retrieve the result using \code{future.get()}.
\end{itemize}

\task{Reflection}
\begin{itemize}
  \item How does \code{wait_for()} differ from \code{get()}?
  \item Is it safe to call \code{get()} after \code{wait_for()} returns
    \code{std::future_status::ready}?
\end{itemize}
